\documentclass[tikz,border=4pt]{standalone}
\usepackage{luatexja-fontspec}
\setmainjfont{Noto Sans CJK JP}
\setsansfont{TeX Gyre Heros}
\setsansjfont{Noto Sans CJK JP}
\renewcommand{\familydefault}{\sfdefault}
\usepackage{amsmath,bm}
% Simulink-style drawing conventions shared by the block diagrams
\usetikzlibrary{arrows.meta,calc,positioning}
\tikzset{
  >={Latex[length=2.1mm,width=1.5mm]},
  blk/.style={draw, line width=0.55pt, fill=white, rectangle, align=center, inner sep=2pt},
  des/.style={blk, fill=blue!6, draw=blue!55!black},
  phy/.style={blk, fill=orange!9, draw=orange!70!black},
  sen/.style={blk, fill=orange!4, draw=orange!70!black, dashed},
  port/.style={draw, line width=0.55pt, rounded corners=2.5mm, minimum width=9mm, minimum height=5mm, fill=white, font=\footnotesize\bfseries},
  sumc/.style={draw, circle, line width=0.55pt, minimum size=5.2mm, inner sep=0pt, fill=white},
  sig/.style={->, line width=0.55pt},
  bus/.style={->, line width=1.5pt},
  sname/.style={font=\scriptsize, inner sep=1pt, fill=white},
  pname/.style={font=\scriptsize, inner sep=1pt, text=black!75},
  title/.style={font=\footnotesize\bfseries, inner sep=0pt},
  dot/.style={circle, fill, inner sep=0pt, minimum size=1.6mm},
}

\begin{document}
\begin{tikzpicture}[x=1cm,y=1cm]
\tikzset{sname/.append style={font=\footnotesize}, pname/.append style={font=\footnotesize}, title/.append style={font=\small\bfseries}}
\def\bw{24mm}\def\bh{20mm}
% ---------------------------------------------------------------- blocks: a staircase from top-left to bottom-right,
% so that every forward signal runs only to the right (the feedback is the one line that has to go back)
\def\yA{0}\def\yB{-1.2}\def\yC{-2.4}\def\yD{-3.6}\def\yE{-4.8}
\node[des, minimum width=\bw, minimum height=\bh] (A) at (3.2,\yA) {};
\node[des, minimum width=\bw, minimum height=\bh] (B) at (7.2,\yB) {};
\node[phy, minimum width=\bw, minimum height=\bh] (C) at (11.2,\yC) {};
\node[sen, minimum width=\bw, minimum height=\bh] (D) at (15.2,\yD) {};
\node[des, minimum width=\bw, minimum height=\bh] (E) at (19.2,\yE) {};
% block texts (title + body), kept inside the block
\def\body#1#2#3{\node[font=\footnotesize, align=center, text width=22mm] at (#1,#2-0.18) {#3};}
\node[title] at (3.2,\yA+0.6) {判断};          \body{3.2}{\yA}{把持前に 1 回\\画像＋プランジャ\\→ 区分 $\hat c$}
\node[title] at (7.2,\yB+0.6) {把持力制御器};  \body{7.2}{\yB}{50\,Hz・積分則\\接触時に\\積分器をリセット}
\node[title] at (11.2,\yC+0.6) {プラント};     \body{11.2}{\yC}{空圧・人工筋肉\\指 3 本・物体\\MuJoCo\,/\,PhysX}
\node[title] at (15.2,\yD+0.6) {センサ};       \body{15.2}{\yD}{内圧・指の曲げ\\布スキン 24 点\\雑音・遅れ}
\node[title] at (19.2,\yE+0.6) {推定器};       \body{19.2}{\yE}{固有感覚で\\把持力 $\hat F$\\接触 $\kappa$}
% ---------------------------------------------------------------- ports
\node[port] (in1) at (0.45,1.9) {1};
\node[pname, anchor=north] at (in1.south) {物体 $o$};
\node[port] (in2) at (0.45,-0.5) {2};
\node[pname, anchor=north] at (in2.south) {取扱い表 $H$};
\node[port] (out1) at (22.4,-1.8) {1};
\node[pname, anchor=north] at (out1.south) {物体の状態};
\node[port] (out2) at (22.4,\yE) {2};
\node[pname, anchor=north] at (out2.south) {推定値 $\hat{\bm y}$};
% ---------------------------------------------------------------- signals (forward: right or down only)
\coordinate (b1) at (1.4,1.9);
\draw[-, line width=0.55pt] (in1.east) -- (b1);
\node[dot] at (b1) {};
\draw[sig] (b1) -- (1.4,0.5) -- (A.west |- 0,0.5);
\draw[sig] (b1) -- (11.2,1.9) -- (C.north);
\node[sname, anchor=south] at (6.3,1.9) {物体 $o$（質量・摩擦・剛性・形）};
\draw[sig] (in2.east) -- (A.west |- in2.east);
\draw[bus] (A.east |- 0,-0.6) -- node[sname, above=2pt] {$F^{*},\,\bar F$} (B.west |- 0,-0.6);
\draw[sig] (B.east |- 0,-1.8) -- node[sname, above=1pt] {$p^{\mathrm{cmd}}$} (C.west |- 0,-1.8);
\draw[bus] (C.east |- 0,-3.0) -- node[sname, above=2pt] {$p,\bm\theta,\bm f_c$} (D.west |- 0,-3.0);
\draw[bus] (D.east |- 0,-4.2) -- node[sname, above=2pt] {$\tilde p,\tilde{\bm\theta},\tilde{\bm s}$} (E.west |- 0,-4.2);
\draw[sig] (C.east |- 0,-1.8) -- node[sname, above=1pt, pos=0.5] {物体の位置・姿勢} (out1.west);
% output and feedback (bus)
\coordinate (b2) at (21.15,\yE);
\draw[bus] (E.east) -- (out2.west);
\node[dot] at (b2) {};
\draw[bus] (b2) -- (21.15,-6.3) -- (5.55,-6.3) -- (5.55,-1.8) -- (B.west |- 0,-1.8);
\node[sname, anchor=north] at (13.35,-6.35) {帰還 $\hat{\bm y}=(\tilde p,\hat F,\kappa)$　力センサの値は使わない};
% legend
\def\lg#1#2#3{\node[#1, minimum width=6mm, minimum height=3.5mm] at (#2,-7.35) {}; \node[font=\footnotesize, anchor=west] at (#2+0.4,-7.35) {#3};}
\lg{des}{0.4}{本研究で設計・学習した部分}
\lg{phy}{5.4}{物理（シミュレーション）}
\lg{sen}{9.9}{センサのモデル（仮定の仕様）}
\draw[sig] (14.8,-7.35) -- (15.6,-7.35); \node[font=\footnotesize, anchor=west] at (15.7,-7.35) {信号};
\draw[bus] (16.8,-7.35) -- (17.6,-7.35); \node[font=\footnotesize, anchor=west] at (17.7,-7.35) {バス};
\node[dot] at (19.0,-7.35) {}; \node[font=\footnotesize, anchor=west] at (19.2,-7.35) {分岐};
\end{tikzpicture}
\end{document}
